% Options for packages loaded elsewhere
% Options for packages loaded elsewhere
\PassOptionsToPackage{unicode}{hyperref}
\PassOptionsToPackage{hyphens}{url}
\PassOptionsToPackage{dvipsnames,svgnames,x11names}{xcolor}
%
\documentclass[
  chinese-hans,
]{ctexart}
\usepackage{xcolor}
\usepackage{amsmath,amssymb}
\setcounter{secnumdepth}{-\maxdimen} % remove section numbering
\usepackage{iftex}
\ifPDFTeX
  \usepackage[T1]{fontenc}
  \usepackage[utf8]{inputenc}
  \usepackage{textcomp} % provide euro and other symbols
\else % if luatex or xetex
  \usepackage{unicode-math} % this also loads fontspec
  \defaultfontfeatures{Scale=MatchLowercase}
  \defaultfontfeatures[\rmfamily]{Ligatures=TeX,Scale=1}
\fi
\usepackage{lmodern}
\ifPDFTeX\else
  % xetex/luatex font selection
\fi
% Use upquote if available, for straight quotes in verbatim environments
\IfFileExists{upquote.sty}{\usepackage{upquote}}{}
\IfFileExists{microtype.sty}{% use microtype if available
  \usepackage[]{microtype}
  \UseMicrotypeSet[protrusion]{basicmath} % disable protrusion for tt fonts
}{}
\makeatletter
\@ifundefined{KOMAClassName}{% if non-KOMA class
  \IfFileExists{parskip.sty}{%
    \usepackage{parskip}
  }{% else
    \setlength{\parindent}{0pt}
    \setlength{\parskip}{6pt plus 2pt minus 1pt}}
}{% if KOMA class
  \KOMAoptions{parskip=half}}
\makeatother
% Make \paragraph and \subparagraph free-standing
\makeatletter
\ifx\paragraph\undefined\else
  \let\oldparagraph\paragraph
  \renewcommand{\paragraph}{
    \@ifstar
      \xxxParagraphStar
      \xxxParagraphNoStar
  }
  \newcommand{\xxxParagraphStar}[1]{\oldparagraph*{#1}\mbox{}}
  \newcommand{\xxxParagraphNoStar}[1]{\oldparagraph{#1}\mbox{}}
\fi
\ifx\subparagraph\undefined\else
  \let\oldsubparagraph\subparagraph
  \renewcommand{\subparagraph}{
    \@ifstar
      \xxxSubParagraphStar
      \xxxSubParagraphNoStar
  }
  \newcommand{\xxxSubParagraphStar}[1]{\oldsubparagraph*{#1}\mbox{}}
  \newcommand{\xxxSubParagraphNoStar}[1]{\oldsubparagraph{#1}\mbox{}}
\fi
\makeatother


\usepackage{longtable,booktabs,array}
\newcounter{none} % for unnumbered tables
\usepackage{calc} % for calculating minipage widths
% Correct order of tables after \paragraph or \subparagraph
\usepackage{etoolbox}
\makeatletter
\patchcmd\longtable{\par}{\if@noskipsec\mbox{}\fi\par}{}{}
\makeatother
% Allow footnotes in longtable head/foot
\IfFileExists{footnotehyper.sty}{\usepackage{footnotehyper}}{\usepackage{footnote}}
\makesavenoteenv{longtable}
\usepackage{graphicx}
\makeatletter
\newsavebox\pandoc@box
\newcommand*\pandocbounded[1]{% scales image to fit in text height/width
  \sbox\pandoc@box{#1}%
  \Gscale@div\@tempa{\textheight}{\dimexpr\ht\pandoc@box+\dp\pandoc@box\relax}%
  \Gscale@div\@tempb{\linewidth}{\wd\pandoc@box}%
  \ifdim\@tempb\p@<\@tempa\p@\let\@tempa\@tempb\fi% select the smaller of both
  \ifdim\@tempa\p@<\p@\scalebox{\@tempa}{\usebox\pandoc@box}%
  \else\usebox{\pandoc@box}%
  \fi%
}
% Set default figure placement to htbp
\def\fps@figure{htbp}
\makeatother



\ifLuaTeX
\usepackage[bidi=basic,shorthands=off,provide=*]{babel}
\else
\usepackage[bidi=default,shorthands=off,provide=*]{babel}
\fi
\ifLuaTeX
  \usepackage{selnolig} % disable illegal ligatures
\fi


\setlength{\emergencystretch}{3em} % prevent overfull lines

\providecommand{\tightlist}{%
  \setlength{\itemsep}{0pt}\setlength{\parskip}{0pt}}



 


\makeatletter
\@ifpackageloaded{tcolorbox}{}{\usepackage[skins,breakable]{tcolorbox}}
\@ifpackageloaded{fontawesome5}{}{\usepackage{fontawesome5}}
\definecolor{quarto-callout-color}{HTML}{909090}
\definecolor{quarto-callout-note-color}{HTML}{0758E5}
\definecolor{quarto-callout-important-color}{HTML}{CC1914}
\definecolor{quarto-callout-warning-color}{HTML}{EB9113}
\definecolor{quarto-callout-tip-color}{HTML}{00A047}
\definecolor{quarto-callout-caution-color}{HTML}{FC5300}
\definecolor{quarto-callout-color-frame}{HTML}{acacac}
\definecolor{quarto-callout-note-color-frame}{HTML}{4582ec}
\definecolor{quarto-callout-important-color-frame}{HTML}{d9534f}
\definecolor{quarto-callout-warning-color-frame}{HTML}{f0ad4e}
\definecolor{quarto-callout-tip-color-frame}{HTML}{02b875}
\definecolor{quarto-callout-caution-color-frame}{HTML}{fd7e14}
\makeatother
\makeatletter
\@ifpackageloaded{caption}{}{\usepackage{caption}}
\AtBeginDocument{%
\ifdefined\contentsname
  \renewcommand*\contentsname{目录}
\else
  \newcommand\contentsname{目录}
\fi
\ifdefined\listfigurename
  \renewcommand*\listfigurename{图索引}
\else
  \newcommand\listfigurename{图索引}
\fi
\ifdefined\listtablename
  \renewcommand*\listtablename{表索引}
\else
  \newcommand\listtablename{表索引}
\fi
\ifdefined\figurename
  \renewcommand*\figurename{图}
\else
  \newcommand\figurename{图}
\fi
\ifdefined\tablename
  \renewcommand*\tablename{表}
\else
  \newcommand\tablename{表}
\fi
}
\@ifpackageloaded{float}{}{\usepackage{float}}
\floatstyle{ruled}
\@ifundefined{c@chapter}{\newfloat{codelisting}{h}{lop}}{\newfloat{codelisting}{h}{lop}[chapter]}
\floatname{codelisting}{列表}
\newcommand*\listoflistings{\listof{codelisting}{列表索引}}
\captionsetup{labelsep=period}
\usepackage{amsthm}
\theoremstyle{definition}
\newtheorem{definition}{定义}[section]
\theoremstyle{remark}
\AtBeginDocument{\renewcommand*{\proofname}{证明}}
\newtheorem*{remark}{注记}
\newtheorem*{solution}{解}
\newtheorem{refremark}{注记}[section]
\newtheorem{refsolution}{解}[section]
\makeatother
\makeatletter
\makeatother
\makeatletter
\@ifpackageloaded{caption}{}{\usepackage{caption}}
\@ifpackageloaded{subcaption}{}{\usepackage{subcaption}}
\makeatother
\usepackage{bookmark}
\IfFileExists{xurl.sty}{\usepackage{xurl}}{} % add URL line breaks if available
\urlstyle{same}
% fallback for those not using the hyperref driver hyperxmp:
\makeatletter
\@ifundefined{xmpquote}{\newcommand{\xmpquote}[1]{#1}}{}
\makeatother
\hypersetup{
  pdftitle={数学分析},
  pdflang={zh-Hans},
  colorlinks=true,
  linkcolor={blue},
  filecolor={Maroon},
  citecolor={Blue},
  urlcolor={Blue},
  pdfcreator={LaTeX via pandoc}}


\title{数学分析}
\author{}
\date{}
\begin{document}
\maketitle


\section{极限}\label{ux6781ux9650}

\subsection{实数}\label{ux5b9eux6570}

\subsubsection{整数与有理数}\label{ux6574ux6570ux4e0eux6709ux7406ux6570}

通常用 \(\mathbb{N}\) 表示自然数集 \[
\mathbb{N} = \{0,1,2,3,\cdots\}
\] 自然数集对加法运算封闭 (定义~\ref{def-closure}) ，对减法运算不封闭

整数集对减法封闭 \[
\mathbb{Z} = \{0,\pm1,\pm2,\pm3,\cdots\}
\] 整数集对乘法运算，对除法运算不封闭

在整数集中添加两个整数的商，就得到了有理数集 \[
\mathbb{Q} =\left\{\frac{p}{q}\middle|p,q \in \mathbb{Z},q \neq 0\right\}
\] 它对加减乘除四则运算封闭，也称 \(\mathbb{Q}\) 为有理数域

有理数具有稠密性（定义~\ref{def-dense-in-R}）

\begin{definition}[]\protect\hypertarget{def-closure}{}\label{def-closure}

\textbf{封闭性}

设 \(S\) 是一个非空集合，\(\circ\) 是定义在 \(S\) 上的一种运算。

若 \(\forall a,b \in S ,都有 a \circ b \in S\)

则称集合 \(S\) 对该运算 \(\circ\) 是\textbf{封闭的}，或称 \(S\)
具有\textbf{封闭性}

\end{definition}

\begin{definition}[]\protect\hypertarget{def-field}{}\label{def-field}

\textbf{域}

设 \(\mathbb{F}\) 是一个包含至少两个元素的集合，在 \(\mathbb{F}\)
上定义了两个运算：加法 \(\oplus\) 和乘法 \(\otimes\)。如果对于
\(\forall a,b,c \in \mathbb{F}\)，满足以下三大类\textbf{域公理}，则称
\(\mathbb{F}\) 为一个\textbf{域}。

\textbf{（一）加法公理}

\begin{enumerate}
\def\labelenumi{\arabic{enumi}.}
\tightlist
\item
  \textbf{封闭性}：\(a \oplus b \in \mathbb{F}\)
\item
  \textbf{交换律}：\(a \oplus b = b \oplus a\)
\item
  \textbf{结合律}：\((a \oplus b)\oplus c = a \oplus(b \oplus c)\)
\item
  \textbf{单位元}：存在元素
  \(0 \in \mathbb{F},\text{ s.t. }a \oplus 0 = a\)
\item
  \textbf{逆元}：存在元素
  \(-a \in \mathbb{F},\text{ s.t. }a \oplus (-a) = 0\)
\end{enumerate}

\textbf{（二）乘法公理}

\begin{enumerate}
\def\labelenumi{\arabic{enumi}.}
\setcounter{enumi}{5}
\tightlist
\item
  \textbf{封闭性}：\(a \otimes b \in \mathbb{F}\)
\item
  \textbf{交换律}：\(a \otimes b = b \otimes a\)
\item
  \textbf{结合律}：\((a \otimes b) \otimes c = a \otimes(b \otimes c)\)
\item
  \textbf{单位元}：存在元素
  \(1 \in \mathbb{F},\text{ s.t. }a \otimes 1 = a\)
\item
  \textbf{逆元}：若 \(a \neq 0\),则存在元素
  \(a^{-1} \in \mathbb{F},\text{ s.t. }a \otimes a^{-1} = 1\)
\end{enumerate}

\textbf{（三）分配律}

\begin{enumerate}
\def\labelenumi{\arabic{enumi}.}
\setcounter{enumi}{10}
\tightlist
\item
  \textbf{乘法对加法的分配律}：\(a \otimes (b \oplus c)= a \otimes b \oplus a \otimes c\)
\end{enumerate}

\end{definition}

\begin{tcolorbox}[enhanced jigsaw, arc=.35mm, bottomrule=.15mm, bottomtitle=1mm, breakable, colback=white, colbacktitle=quarto-callout-note-color!10!white, colframe=quarto-callout-note-color-frame, coltitle=black, left=2mm, leftrule=.75mm, opacityback=0, opacitybacktitle=0.6, rightrule=.15mm, title=\textcolor{quarto-callout-note-color}{\faInfo}\hspace{0.5em}{注记}, titlerule=0mm, toprule=.15mm, toptitle=1mm]

这里的加法 \(\oplus\) 和乘法 \(\otimes\) 不一定是经典认知上的 \(+\) 和
\(\times\)
，而也可以是其他运算，加法单位元和乘法单位元也不一定是0和1，比如：对于正实数集
\(\mathbb{R}_+\),在其中定义加法及乘法运算为 \[
a \oplus b = ab\quad(a,b \in \mathbb{R}_+) \\
a \otimes b = a^b\quad(a,b \in \mathbb{R}_+)
\] 可以验证这种情况下， \(\mathbb{R}_+\)
是一个域，其加法单位元为1，乘法单位元也为1。

\end{tcolorbox}

\begin{definition}[]\protect\hypertarget{def-dense-in-R}{}\label{def-dense-in-R}

\textbf{（狭义）稠密性} 设 \(S\) 是实数集 \(\mathbb{R}\) 的一个子集。若
\[
\forall a,b \in S,a\neq b,\exists c \in S,\text{s.t.} a<c<b
\] 则称集合 \(S\) 在 \(\mathbb{R}\) 中\textbf{稠密}，或者称 \(S\)
具有\textbf{稠密性}

\end{definition}




\end{document}
